% =====================================================================
%  RESUMEN Y ABSTRACT — VERSIÓN CORREGIDA
%  Ajuste de cohesión y claridad en los últimos tres párrafos.
%  El párrafo final refleja de forma concreta lo que se cumplió a lo
%  largo del trabajo. Redacción sencilla, coherente y sin tecnicismos
%  innecesarios para que cualquier lector pueda seguirla.
% =====================================================================

\frontsection{RESUMEN}
La anotación de videos en Lengua de Señas Ecuatoriana (LSEc) exige herramientas que respeten el criterio lingüístico del investigador y, al mismo tiempo, permitan aprovechar modelos modernos de Inteligencia Artificial. ELAN, una de las herramientas más utilizadas para la anotación multimodal, facilita segmentar y etiquetar eventos sobre una línea de tiempo. No obstante, su arquitectura basada en Java y su mecanismo de extensión dificultan la conexión directa con servicios de inferencia desarrollados en entornos actuales, principalmente Python, con dependencias especializadas. Esto limita la incorporación de procesos automáticos, como la segmentación temporal y la clasificación de glosas, al flujo habitual de anotación.

Para resolver esta limitación, en este trabajo se desarrolló un middleware que funciona como un intermediario entre ELAN y los modelos de Inteligencia Artificial, sin necesidad de modificar a fondo la herramienta. La idea central es sencilla: ELAN actúa como cliente y envía la solicitud de análisis, mientras que el middleware recibe esa solicitud, verifica que los datos sean correctos, ubica el modelo que debe usarse y coordina su ejecución en un entorno aislado. Finalmente, el resultado regresa a ELAN convertido en segmentos de tiempo que el investigador puede aceptar, ajustar o eliminar. De este modo, la revisión final siempre queda en manos del especialista.

El desarrollo se llevó a cabo por etapas, avanzando paso a paso y comprobando cada avance antes de continuar con el siguiente. Primero se estudió cómo funciona ELAN por dentro para entender la forma correcta de conectarle un proceso externo. Con esa base se definió una organización en tres niveles: ELAN como herramienta de anotación, el middleware como coordinador local y los modelos como servicios independientes que procesan los videos. A partir de ahí se fueron sumando capacidades: la comunicación básica, la instalación de modelos, su ejecución controlada, el seguimiento de cada tarea y, por último, la conexión visual dentro de ELAN.

Los resultados confirmaron que la integración funciona y es estable. El sistema permitió instalar modelos, comprobar que estén disponibles, ejecutar análisis sobre los videos, informar en qué estado se encuentra cada tarea y distinguir con claridad entre un error real y una demora por tiempo agotado. Las pruebas realizadas mostraron que el middleware responde de forma ordenada tanto en situaciones normales como ante problemas frecuentes, por ejemplo un video que no existe, un modelo desactivado o un contenedor detenido, sin que ello afecte a ELAN. En conjunto, el trabajo entregó una solución local, ordenada y ampliable que acerca los modelos de Inteligencia Artificial al trabajo diario de anotación, mantiene los datos dentro del equipo del investigador y deja el camino preparado para incorporar nuevos modelos en el futuro sin tener que construir otra herramienta.

\textbf{PALABRAS CLAVE:} ELAN, middleware, orquestación de servicios, Inteligencia Artificial, Lengua de Señas Ecuatoriana.

\frontsection{ABSTRACT}
Video annotation in Ecuadorian Sign Language (LSEc) requires tools that preserve the researcher's linguistic judgment while making it possible to use modern Artificial Intelligence models. ELAN is widely used for multimodal annotation because it supports segmenting and labeling events on a timeline. However, its Java-based architecture and extension mechanism make direct integration difficult with inference services developed in current environments, especially Python, with specialized dependencies. This limits bringing automatic processes, such as temporal segmentation and gloss classification, into the regular annotation workflow.

To address this limitation, this work developed a middleware that acts as an intermediary between ELAN and Artificial Intelligence models, without having to deeply modify the tool. The core idea is simple: ELAN works as a client and sends the analysis request, while the middleware receives that request, checks that the data is correct, locates the model to be used, and coordinates its execution in an isolated environment. Finally, the result returns to ELAN as time segments that the researcher can accept, adjust, or delete. In this way, the final review always remains in the hands of the specialist.

The development was carried out in stages, moving step by step and verifying each advance before continuing with the next one. First, the inner workings of ELAN were studied to understand the correct way to connect an external process to it. On that basis, a three-level organization was defined: ELAN as the annotation tool, the middleware as the local coordinator, and the models as independent services that process the videos. From there, capabilities were added progressively: basic communication, model installation, their controlled execution, the tracking of each task, and, finally, the visual connection within ELAN.

The results confirmed that the integration works and is stable. The system made it possible to install models, verify that they are available, run analyses on the videos, report the state of each task, and clearly distinguish a real error from a delay due to a timeout. The tests showed that the middleware responds in an orderly way both in normal situations and when facing common problems, such as a video that does not exist, a disabled model, or a stopped container, without affecting ELAN. Overall, the work delivered a local, orderly, and extensible solution that brings Artificial Intelligence models closer to the daily annotation work, keeps the data inside the researcher's own machine, and leaves the path prepared to incorporate new models in the future without having to build another tool.

\textbf{KEYWORDS:} ELAN, middleware, service orchestration, Artificial Intelligence, Ecuadorian Sign Language.
